%-------------------------------------------------------------------------------
\section{Introduction}
%-------------------------------------------------------------------------------

%o evaluate the performance of a research prototype and compare different prototypes,
%researchers often rely on well-established benchmarks, such as TPC-C~\cite{tpccurl} and YCSB~\cite{ycsburl,Cooper2010YCSB}.
Through the emphasis on artifact evaluation, our community has greatly
improved the quality of artifacts and simplified the process of reproducing and comparing to prior results.

However, we observe one long-lasting question persists despite the
improvement of reproducibility: which settings (i.e., combinations of parameter values) should we use
during our experiments? Popular benchmarks like TPC-C~\cite{tpccurl} and YCSB~\cite{ycsburl,Cooper2010YCSB}
and the systems to evaluate often have multiple parameters; as shown in prior works~\cite{Harding2017EvalDistributed,Wu2017EvalMem,Lim2017Cicada,Lu2020Aria}, 
it often happens that one system can outperform another under one setting
but not so under a different setting.
So far our artifact evaluation only requires reproducing results under the
same settings as reported in the corresponding publication, but does not
require or recommend any settings to use.

In general, there exist two philosophies to answer this question. First, we can focus on one or a few
representative settings, usually extracted from studies of production systems. 
%This approach
%has the advantage that we know any improvement under these settings will
%have a real-world impact, but if we insist that any works must follow these
%settings, it may discourage works exploring new problems. 
Second, we can
perform an extensive evaluation to cover all/many combinations of parameter values.
Extensive evaluation has the obvious benefit that it can provide a comprehensive
picture of system performance, which can avoid potential bias introduced by
parameter selection, but it may cost a large amount of time and resources.

%which will allow readers to have a full picture. However, if the
%benchmark includes many parameters or many possible values, this
%approach will be time-consuming. 
%Note that these two philosophies are
%not completely orthogonal to each other since a number of studies find
%industrial companies can have a variety of workloads, which call for extensive
%evaluation as well.

Which approach is more appropriate perhaps requires a thorough discussion
in our community, but this paper tries to advocate extensive evaluation
due to the following two observations:
first, a number of studies have shown that industrial companies
can have a variety of workloads~\cite{Atikoglu2012WAL,Yang2020Twitter,Nishtala2013Scaling,Berg2020CacheLib}.
For example, Facebook finds ``object sizes spanning
more than seven orders of magnitude'' for certain applications~\cite{Berg2020CacheLib}.
%In this case, 
%optimizing for the
%most common workload is certainly beneficial, but ignoring or incurring
%significant slowdown for less common workloads is not ideal~\cite{}.
Such observation naturally calls for
an extensive evaluation to reveal the full picture of system performance under different settings.
Second, we observe a later research work often needs to compare with a prior work under a
new setting that was not tested by the prior work. Our experience shows
that this often incurs many problems, usually due to bugs or performance
issues that only manifest under the new settings (Section~\ref{sec:cases}); a number of works choose
to re-implement prior works (e.g., Janus~\cite{Mu2016Consolidating} re-implements TAPIR~\cite{Zhang15Building}; Aria~\cite{Lu2020Aria} re-implements
Calvin~\cite{Thomson2012Calvin}, etc)---this is hurting the purpose of artifact evaluation, i.e., ``easy to reuse, facilitating further research''~\cite{osdi-repro} to
some degree. An extensive evaluation can help a researcher
to identify and fix such problems.

%However, we observe many works in our community are neither evaluated under a
%representative setting, often due to the fact that we don't even know the
%representative settings, nor evaluated extensively, probably due to time or
%resource constraints: this is certainly not ideal. Our emphasis on artifact
%evaluation does not address this problem since it neither restricts parameter
%values nor requires experiments beyond settings used in the paper. 

This paper studies the feasibility and benefits of extensive evaluation.
It tries to answer the following questions: 1)  Can extensive evaluation help us to improve the quality of artifacts
and allow a more comprehensive performance comparison?
2) In terms of time and resources, is it feasible
to carry out extensive evaluation? And if not, how can we address that?


To answer these questions, we focus on TPC-C and YCSB, two of the most popular
benchmarks in our community, because there is a general consensus about what
parameters can be tuned for these two benchmarks.
We selected 10 research prototypes evaluated
under these two benchmarks. In addition to reproducing their results as
reported in their original papers, we further run settings that are not used
in these papers: when we see failed tests or significant performance degradation,
we try our best effort to diagnose and fix the problems. From such experience,
we try to answer the above questions in terms of both system correctness and
performance.

\vspace{-.1in}
\paragraph{Correctness.} In our experiments, it is not uncommon for a
research prototypes to stop working under untested settings. However, since even
mature commercial softwares are not free of bugs, it is perhaps unrealistic
to expect a research prototype to be stable and reliable under all conditions,
so the question we want to investigate here is whether it is possible to improve
the quality of artifact with a reasonable amount of effort.

We tried to diagnose and fix the problems we met and observe the problems can
be categoried into two types: the first type deterministically blocks an experiment
under a certain setting. They are usually caused by hard-coded configurations.
For example, a prototype may preallocate a fixed amount of memory, which will
not work if we test the prototype with more data. Naturally, this type of problems
are more likely to happen when we push the system to the extreme (e.g., a lot of
data or a lot of nodes). Diagnosing this type of problems could take a while for
us, since the error messages do not always directly point to the root cause, but should be
easier for the authors who are familiar with the system; fixing this type of problems
usually only requires changing a few lines of code. The second type can non-deterministically
cause the prototype to crash during certain experiments. While we try our best to
diagnose this type, we only succeed in a subset of them. This type of problems
can happen under any settings, which means identifying them will need a large
amount of effort, not to mention diagnosing and fixing them.

From such experience, we conclude that, while expecting a research prototype
to be almost bug-free is unrealistic, fixing the first type of problems is both
beneficial and feasible: it allows a future research work to compare with a prior
work under a different setting, without requiring the additional effort of
fixing the code of prior work; these types of problems can usually be triggered
by running extreme settings and can usually be fixed easily, which means the
authors do not need to spend a significant amount of effort to test different
settings and diagnose the problems.

\vspace{-.1in}
\paragraph{Performance.}
When comparing system performance under a wide range of settings,
we observe extensive evaluation can bring two benefits:

First, extensive evaluation has revealed design issues that
    cause performance degradation in previously untested settings. For example, DrTM~\cite{Wei2015FIT}
    and Silo experience performance degradation when loaded with
    more data, since their tree indexes grow larger. This problem could be addressed
    by partitioning data~\cite{Kalia2014HERD,chang06bigtable,Decandia07Dynamo,Lee2021ShardManager}.
    %, though this requires quite significant changes and thus
    %is out of the scope of this work. 
    Revealing and addressing these issues can avoid
    the problem that a later work has to ``optimize'' or ``re-implement'' a prior work for
    a fair comparison under a new setting. Furthermore, we also observe
    performance degradation that we do not know how to address, which
    could be future research topics as well.
    
Second, extensive evaluation allows us to perform a more comprehensive comparison
    of system performance under different settings. Perhaps
    not surprisingly, we find prior works sometimes report
    settings they excel at: while these settings may be important,
    it may not be a good idea to completely ignore other settings,
    as discussed above.

The cost, however, is significant: it tooks us thousands of machine hours to test
one prototype, despite that we carefully chose the parameters to tune and
 the parameter values to use. For systems with more parameters, the cost
could get out of control.

To address this challenge, we investigate the following question: is it possible
to identify all such performance patterns by running a subset of settings?

......



\iffalse
\vspace{-.1in}
\paragraph{Sampling and Regression to Reduce Experiment Time.}
Our experience shows that, perhaps not surprisingly, extensive
evaluation can indeed be quite  time- and resource-consuming, in particular
when evaluating distributed systems.
For experiments with 4-5 parameters, we have to carefully choose parameter
values to make total experiment time manageable.
If we need to add a few more parameters or
possible values, the experiment time will quickly get out of control. 

To address this problem,
we have tried the classic sampling and regression approach: we sample a subset of
settings, measure the performance under these
settings, and use the result to predict performance under other settings.
We make two observations. First, non-linear regression methods like 
multilayer perceptron (MLP)~\cite{MLP1998}
works quite well. Given enough samples, MLP can achieve a normalized root-mean-square deviation (RMSE) value of 0.1 or lower for most systems.
Second, for systems with less predictable performance, we find they often
have performance anomalies under certain settings.
These phenomena are perhaps not surprising because unpredictability indicates
that system performance can have an abrupt change due to a small change of setting:
this is undesired for almost any real-world system.

Motivated by these two observations, we suggest the following approach to determine
what settings to run: we can gradually measure more settings, until we find 
the normalized RMSE of regression is stable. In this case, we can be
confident that we have run enough samples to reveal the system's performance
characteristics. Furthermore, if the normalized RMSE is high,
we should investigate why. 
%Regression can help us reveal the settings whose
%performance often deviate from prediction, which are the targets we should investigate.
%Note that even for systems with well predictable performance, there might
%still be a few such abnormal settings.

\vspace{-.1in}
\paragraph{Tiled Heat Map to Visualize High-Dimensional Data.}
Since the benchmarks include multiple
parameters, the evaluation results are high-dimensional data that cannot
be directly mapped to a two-dimensional figure. If we only present a subset
of results, we will need to go back to the question of which settings
to use---this is exactly what we try to avoid by
extensive evaluation. 
%There are a lot of works about how to
%visualize high-dimensional data, but most focus on finding the
%structures within high-dimensional data, which is different from
%our goal of comprehensive comparison.

To address this problem,
we design and develop a tool to visualize high-dimensional data on
a figure called Tiled Heat Map: for $n$ dimensional
data, this tool first chooses $n-2$ dimensions and groups data with the same
values of these two $n-2$ dimensions; each group only has two
dimensions and thus can be drawn in a 2D figure, called a ``tile'';
then this tool treats each tile as a point and repeats this procedure for
the remaining $n-2$ dimensions, till it can cover all dimensions.
To determine which dimensions to choose in each round, we further
propose a metric to evaluate how ``smooth'' the result
figure is by computing the correlation of the distance and
the difference between all pairs of points: our tool will iterate
possible choices and choose the one with the highest smoothness.
In our study, we find the Tiled Heat Map can well present 4D-5D data to
show the trend of how parameter values affect system
performance.


\vspace{-.1in}
\paragraph{Benefits of Extensive Evaluation.}
Our experience shows that extensive evaluation can improve
the quality of artifact and allow a more comprehensive
performance analysis:

\begin{itemize} [leftmargin=*]

    \vspace{-.05in}
    \item Extensive evaluation has revealed multiple correctness
    or configuration issues that block extensive evaluation.
    Some of them are caused by hard-coded configurations (e.g., fixed memory allocation by Silo~\cite{Tu2013Silo}, Cicada~\cite{Lim2017Cicada}, DrTM~\cite{Wei2015FIT}, and GAM~\cite{Cai2018GAM}),
    which can easily be fixed by the authors but can take much longer
    for others to figure out. Some of them are bugs that are more
    likely to manifest under new settings. Revealing and fixing these
    issues is exactly one purpose of artifact evaluation and that's why
    we further advocate extensive evaluation.
    Furthermore, we find a fundamental solution
    to some of these problems (e.g., adaptive configuration rather than manual
    tuning) is beyond simple engineering effort and perhaps is a research topic by itself.
    
    \vspace{-.05in}
    \item Extensive evaluation has also revealed design issues that
    cause performance degradation in previously untested settings. For example, DrTM~\cite{Wei2015FIT}
    and Silo experience performance degradation when loaded with
    more data, since their tree indexes grow larger. This problem could be addressed
    by partitioning data~\cite{Kalia2014HERD,chang06bigtable,Decandia07Dynamo,Lee2021ShardManager}.
    %, though this requires quite significant changes and thus
    %is out of the scope of this work. 
    Revealing and addressing these issues can avoid
    the problem that a later work has to ``optimize'' or ``re-implement'' a prior work for
    a fair comparison under a new setting. Furthermore, we also observe
    performance degradation that we do not know how to address, which
    could be future research topics as well.
    
    \vspace{-.05in}
    \item Extensive evaluation allows us to perform a more comprehensive comparison
    of system performance under different settings. Perhaps
    not surprisingly, we find prior works sometimes report
    settings they excel at: while these settings may be important,
    it may not be a good idea to completely ignore other settings,
    as discussed above.
        
\end{itemize}
\fi